\section{An exact classification for congruence-filtered prime sums}
\label{sec:filters}
Fix an even integer \(q\geq2\) and a nonempty subset
\(R\subseteq(\mathbb Z/q\mathbb Z)^\times\).
For each even \(N\), let \(G_N(q,R)\) join distinct integers
\(i,j\in[1,N]\) precisely when \(i+j\) is prime and
\(i+j\bmod q\in R\).
The modulus and the set \(R\) are fixed before \(N\) tends to infinity.
Put
\[
 \sigma_R=\frac{|R|}{\phiG(q)},\qquad
 g_R=\gcd\bigl(q,\{r-r':r,r'\in R\}\bigr).
 \tag{33}
\]
For a singleton set the difference list consists of zero,
so \(g_R=q\). Representatives do not affect the gcd.
Every element of \(R\) is odd, so \(g_R\) is even.

Write
\[
 \rho_R(b)=\frac1{|R|}\sum_{r\in R}\e(br/q),\qquad
 p_q=\min\{p:p\ {\rm prime},\ p\nmid q\},
\]
and, when \(g_R=2\), define the positive explicit spectral gap
\[
 \delta_R=1-\max\left\{\frac1{p_q-1},
       \max_{\substack{b\bmod q\\b\ne0,\ q/2}}|\rho_R(b)|
                  \right\}>0.                              \tag{34}
\]
An empty inner maximum is zero.

\begin{theorem}[The finite obstruction is complete]
\label{thm:filterclassification}
For fixed \(q,R\), the following are equivalent:
\begin{enumerate}
\item \(g_R=2\).
\item \(G_N(q,R)\) is connected for every sufficiently large even \(N\).
\item \(G_N(q,R)\) is Hamiltonian for every sufficiently large even \(N\).
\end{enumerate}
When these hold, \(G_N(q,R)\) has a spanning regular subgraph
of degree
\[
 r_N=(1-o(1))\sigma_R\frac N{\log N}
\]
with nontrivial spectral ratio at most \(1-\delta_R+o(1)\).
For every fixed \(0<\kappa<\delta_R/4\) it contains at least
\[
 \left\lfloor\kappa\sigma_R\frac N{\log N}\right\rfloor
 \tag{35}
\]
edge-disjoint Hamilton cycles.
It remains Hamiltonian after arbitrary edge deletions of
maximum degree at most
\[
 \frac{\delta_R^2}{1000}\sigma_R\frac N{\log N}.              \tag{36}
\]
If \(g_R>2\), the graph is disconnected for every \(N\geq g_R\).
\end{theorem}

\begin{lemma}[The complete signed Fourier model]
\label{lem:filtermodel}
Let \(\Gamma=\mathbb Q/\mathbb Z\).
The compact self-adjoint model for the normalized weight
\[
 A_N^R(i,j)=\frac{\phiG(q)}{|R|N}
                     \Lambda(i+j)\one_{\,i+j\bmod q\in R}
\]
acts on the orthonormal basis \((e_\gamma)_{\gamma\in\Gamma}\) by
\[
 C_Re_\gamma=w_R(\gamma)e_{-\gamma}.
 \tag{37}
\]
The coefficient is zero unless
\[
 \gamma=b/q+a/c,\qquad
 (c,q)=1,\quad c\ {\rm squarefree},\quad(a,c)=1,
\]
where \(c=1,a=0\) is included.
In the nonzero cases its exact value is
\[
 w_R(\gamma)=\rho_R(b)\frac{\mu(c)}{\phiG(c)}.                \tag{38}
\]
If \(g_R=2\), the only eigenvalues of modulus 1 are the simple
constant and parity eigenvalues \(1,-1\).
Every other eigenvalue has modulus at most \(1-\delta_R\).
The sampled operator has the same spectral bounds with an \(o(1)\)
error.
\end{lemma}
\begin{proof}
One can view (38) as a finite character computation, with no
hypothesis about unproved prime correlations.
The limiting prime residue law is the product of the normalized
uniform unit laws at each prime power, conditioned on its residue
modulo \(q\) lying in \(R\).
At primes dividing \(q\), additive characters of conductor
larger than the corresponding prime-power factor of \(q\)
average to zero over the uniform lifts.
At any prime not dividing \(q\), a primitive additive character
averages to \(-1/(p-1)\) at conductor \(p\), and to zero at
conductor \(p^k\), \(k\geq2\).
The Chinese remainder theorem makes these latter coordinates
independent of the residue modulo \(q\).
The remaining finite average at \(q\) is exactly \(\rho_R(b)\).
This proves every zero coefficient and the entire signed formula
(38). In particular
\(w_R(-\gamma)=\overline{w_R(\gamma)}\), so (37) is self-adjoint;
its two-dimensional reflection blocks have eigenvalues
\(\pm|w_R(\gamma)|\).
Conductor tails tend to zero, since \(q\) is fixed and
\(\phiG(c)\geq\sqrt{c/2}\), so the model is compact.

Here are details linking this exact model to the sampled matrix.
Let \(C\) be the unfiltered reflection model with coefficients
\(\mu(k)/\phiG(k)\) at reduced conductor \(k\).
For \(t=b/q\), let \(U_t e_\gamma=e_{\gamma+t}\).
The finite Fourier expansion of the residue mask gives the
operator identity
\[
 C_R=\frac{\phiG(q)}{|R|q}
       \sum_{b\bmod q}
       \left(\sum_{r\in R}\e(-br/q)\right)U_{b/q}CU_{b/q}.
 \tag{39}
\]
Indeed the coefficient at \(\gamma\) in its reflection is the
finite masked average of \(\e(\gamma t)\) under the prime
residue law just computed.
Notice the same shift on both sides of \(C\), rather than an
adjoint conjugation: this is a sum kernel.

Truncate \(C\) to reduced conductors \(k\leq U\).
The operator tail is at most \(\sqrt{2/U}\).
The finite sum (39) therefore has tail \(O_{q,R}(U^{-1/2})\).
All frequencies in its truncated range and domain have
denominator at most \(qU\).
Their sampled Gram matrix differs from the identity by
\(O_q(U^4/N)\), by geometric summation as in
Lemma~\ref{lem:integerspectrum}.
On the integer samples the same finite Fourier expansion
takes \(A_N\) to \(A_N^R\): multiplication by
\(\e(b(i+j)/q)\) is a diagonal unitary on each side.
Lemma~\ref{lem:integertransfer}, with \(U\) a sufficiently
large fixed power of \(\log N\), bounds the unfiltered
operator error by any prescribed negative power of \(\log N\).
The finite mask expansion multiplies this bound only by a
constant depending on \(q,R\).
The Gram bound and the truncated model tail tend to zero.
Polar orthogonalization and Weyl's inequality therefore
transfer the spectrum of \(C_R\), with appended zeros, to
\(A_N^R\). This also justifies the model without appealing to
a pointwise infinite Ramanujan expansion.

It remains to locate its unit-modulus coefficients.
A finite average of unit complex numbers has modulus 1
exactly when all its terms agree. Thus
\[
 |\rho_R(b)|=1
 \quad\Longleftrightarrow\quad
 q\mid b(r-r')\ \hbox{for all }r,r'\in R
 \quad\Longleftrightarrow\quad
 b\equiv kq/g_R\pmod q
\]
for some integer \(k\).
When \(g_R=2\), these are just \(b=0,q/2\), with averages
\(1,-1\).
For \(c>1\), (38) has modulus at most \(1/(p_q-1)\).
Every remaining coefficient is bounded by (34).
This proves all the assertions.
\end{proof}

\begin{proof}[Proof of Theorem~\ref{thm:filterclassification}]
First suppose \(g_R>2\), and fix \(r_0\in R\).
Every edge satisfies
\[
 i+j\equiv r_0\pmod {g_R}.
\]
Since \(g_R\) is even and \(r_0\) odd, the involution
\(a\mapsto r_0-a\) on residues modulo \(g_R\) has no fixed point.
Its \(g_R/2\) distinct pairs give disjoint vertex sets with
no edge between them. All are nonempty once \(N\geq g_R\).
There are at least two such sets, so the graph is disconnected.

Now suppose \(g_R=2\).
Remove proper prime powers and diagonal entries from \(A_N^R\),
and denote the resulting true prime weight matrix by \(Q_N^R\).
The removed maximum row sum is
\(O_{q,R}(N^{-1/2}(\log N)^2)\).
The prime number theorem in the fixed reduced classes modulo
\(q\), with arbitrary logarithmic accuracy, gives uniformly
\[
 \|Q_N^R\one-\one\|_\infty\ll_{q,R,B}(\log N)^{-B}
 \quad\hbox{for every fixed }B.
\]
Indeed the exact raw row is the normalized difference of the
restricted Chebyshev function at \(i+N\) and \(i\).
The normalized spectrum has the two simple exceptions and
gap \(\delta_R-o(1)\), by Lemma~\ref{lem:filtermodel}.
All entries lie between opposite parities, since \(q\) is even
and \(R\) contains only units.
Lemma~\ref{lem:scaling} therefore gives a true-support
stochastic \(P_N^R\), with
\[
 \|P_N^R-Q_N^R\|=o(1),\qquad
 \|P_N^R|_{\{\one,s\}^\perp}\|\leq1-\delta_R+o(1).
\]
Put \(D=\sigma_R N/\log N\), and let \(B_N^R\) be the
ordinary adjacency.
The row estimate in Lemma~\ref{lem:unweightedtransfer}, with
the fixed normalization factor \(1/\sigma_R\), gives
\[
 \max_i|\deg(i)/D-1|=o(1),\qquad
 \|B_N^R/D-P_N^R\|=o(1).                                    \tag{40}
\]
The Laplacian calculation proving (28) now gives the cut
constant \(g=\delta_R/3\) for all sufficiently large \(N\).
Lemma~\ref{lem:regularfactor} gives an
\(r_N=(1-o(1))D\) factor, discarding only \(o(D)\)
edges at each vertex.
Equation (40) transfers the nontrivial spectral bound to this
factor. Its polynomial-order degree dominates the fixed
polylogarithmic threshold in (2), giving a Hamilton cycle.
The iterative estimate (31), with initial gap \(\delta_R\)
in place of \(1/2\), proves (35).
Hamiltonicity implies connectedness, completing the equivalence.

For completeness we verify the uniform deletion constant (36).
Write \(\delta=\delta_R\) and \(\rho=\delta^2/1000\).
Choose \(N\) so large that the initial normalized degree defect
is at most \(\rho\), the initial second eigenvalue is at most
\((1-\delta+\rho)D\), and the cut constant is \(\delta/3\).
After deletion of maximum degree \(\rho D\), the degree defect
is at most \(2\rho\), and the cut constant is at least
\(\delta/3-\rho\geq\delta/4\).
The regular-factor lemma gives
\[
 r=\left\lfloor(1-16\rho/\delta)D\right\rfloor
                     \geq(1-17\rho/\delta)D
\]
for sufficiently large \(D\).
The total discarded degree from the original graph is at most
\(18\rho D/\delta\). Hence its factor satisfies
\[
 r-\lambda_2(H)
 \geq\left(\delta-\rho-\frac{35\rho}{\delta}\right)D
 \geq \frac{\delta}{2}D.
\]
The factor is regular, bipartite, and of degree comparable to
\(D\), so (2) applies again. This proves (36).
\end{proof}

\begin{corollary}[Arbitrarily small fixed prime density suffices]
\label{cor:sparseprimeclasses}
For every fixed even \(q\geq4\), permit only prime sums
congruent to \(1\) or \(-1\) modulo \(q\).
The resulting graphs are Hamiltonian at every sufficiently
large even order.
The permitted primes have relative density \(2/\phiG(q)\)
among all primes. This can be made arbitrarily small by
choosing \(q\) first and then allowing \(N\) to tend to infinity.
\end{corollary}
\begin{proof}
The difference of the two residues is 2, so \(g_R=2\).
The classification applies, and the fixed-progression prime
number theorem gives the stated density.
There is no claim uniform in a growing modulus.
\end{proof}

\begin{remark}[What the criterion detects]
A single reduced residue class modulo even \(q>2\) fails the
criterion and forces disconnection. Two opposite classes pass.
Thus density alone is not the deciding invariant.
The decisive datum is the subgroup generated by differences
of allowed prime residues, and the retained Fourier phases
compute its obstruction exactly.
\end{remark}

\begin{remark}[Equal density, different outcomes, and a nonreal phase]
For \(q=8\), all three rows below retain exactly half of
the primes asymptotically.
\[
 \begin{array}{c|c|c|c}
 R&\sigma_R&g_R&\hbox{eventual outcome}\\ \hline
 \{1,3\}&1/2&2&\hbox{Hamiltonian}\\
 \{1,5\}&1/2&4&\hbox{disconnected}\\
 \{1,7\}&1/2&2&\hbox{Hamiltonian}
 \end{array}
\]
For the first row, the exact coefficient at \(b=1\) is
\[
 \rho_{\{1,3\}}(1)
 =\tfrac12\bigl(\e(1/8)+\e(3/8)\bigr)
 =\mathrm{i}/\sqrt2 .
\]
Its reflection pair has eigenvalues \(\pm1/\sqrt2\);
the imaginary phase is retained before the eigenvalues are
taken. Here \(\delta_R=1-1/\sqrt2\) and the counting law is
\[
 h(G_N(8,\{1,3\}))
   =\left(\frac{N}{2e\log N}\right)^N\exp(o(N)).
\]
The middle row instead has no Hamilton cycle for any
\(N\geq4\), because its two residue-pair components cannot
communicate.
\end{remark}
